
%\needspace{4\baselineskip}
\color{uniblue}{\subsection{ Автоматический ввод резерва (АВР)}}
\color{black}

\begin{enumerate}

%%%%%%%%%%%%%%%%%%%%%%%%%%%%%%%%%%%%%%%%%%%%%%%%%%%%%%%%%%%%%%%%%%%%%%  ОБЩИЕ СВЕДЕНИЯ %%%%%%%%%%%%%%%%%%%%%%%%%%%%%%%%%%%%%%%%%%%%%%%%%%
\item Описание функции АВР

\begin{enumerate}

\item
В устройстве автоматический ввод резерва представлен функцией «АВР». Функция  предназначена  для  восстановления  питания  потребителей  путем  автоматического присоединения резервного источника.

\item
Данный алгоритм может быть применен для рабочего или резервного ввода, что определяется уставкой программного переключателя SGF2 «Тип\_ввода».

\item
В  нормальном  состоянии  выключатель  ВВ  рабочего  ввода  включен,  а  секционный  выключатель отключен.  При  первом  включении  выключателя  рабочего  ввода  и  отсутствии  сигнала  запрета  АВР, таймером Т2, начинает набираться время готовности АВР подготавливая схему АВР к работе.

\item
На  устройстве, установленном  на  рабочем  вводе  уставка  должна  быть  выставлена  в  положение «Рабочий». 

\item
В  качестве критерия включения  АВР резервного ВВ  или  СВ  необходимо  одновременное  наличие следующих условий:
\begin{itemize} 
    \item наличие разрешения АВР от СВ (или резервного ввода) «Разрешение АВР»; 
    \item сигнал аварийного отключения выключателя из логики КСВ рабочего ВВ; 
    \item с контролем или без контроля отсутствия напряжения на секции, определяемые уставкой программного переключателя SGF3 «Реж\_КОН». 
\end{itemize}

\item
Ввод АВР определяется уставкой таймера Т3 с нерегулируемой выдержкой времени равной 10 с. 

\item
При наличии готовности и условий ввода АВР по истечении регулируемой выдержки времени таймера Т1 подается команда «АВР: Вкл. рез.ВВ или СВ» в терминал на включения СВ (или выключателя резервного ввода).  Таймер  Т6  (1 с)  применяется  для  расширения  импульсного  сигнала  включения  обеспечивая надежное включение выключателя СВ или резервного ввода. 

\item
При срабатывании АВР происходит сброс таймера Т3 и готовности АВР в данном режиме работы схемы и осуществляется при наличии одного из условий:
\begin{itemize}  
    \item наличие сигнала «АВР: Вкл. рез.ВВ или СВ»; 
    \item активный сигнал из логики запрета АВР «Запрет АВР»; 
    \item функция АВР выведена из работы.
\end{itemize}

\item
При аварийном отключении выключателя рабочего ввода формируется сигнал «КСВ: В аварийно отключен», которым производится пуск АВР рабочего ввода. После набора выдержки времени готовности подается команда «АВР: Вкл. рез. ВВ или СВ» на включение СВ (или выключателя резервного ввода), при этом срабатывает выходное реле терминала ВВ. 

\item
На устройстве, установленном на СВ (резервном вводе) уставка SGF2 «Тип\_ввода» должна быть выставлена в положение «Резервный». 

\item
В нормальном состоянии выключатель ВВ рабочего ввода включен, а выключатель резервного ввода (секционный выключатель СВ) -- отключен.

\item 
При отключенном положении СВ (или выключатель резервного ввода) «В: Отключен» и отсутствии сигнала запрета АВР таймером Т2 начинает набираться время готовности АВР подготавливая схему АВР к работе.

\item
Разрешение АВР в СВ (или в выключателе резервного ввода) формируется при поступлении сигналов наличия  напряжения  на  вводном  выключателе  или  на  секции  шин.  Контроль  встречного  напряжения определяется уставкой SGF4 «Контр\_встреч\_U».

\item
В  качестве  критерия  включения  АВР  СВ  (или  выключателя  резервного  ввода)  необходимо 
одновременное наличие следующих условий:
\begin{itemize} 
    \item наличие готовности АВР; 
    \item активный дискретный сигнал «Включение АВР» от основного ввода; 
    \item с контролем или без контроля отсутствия напряжения на секции, определяемые уставкой SGF3 «Реж\_КОН». 
\end{itemize}

\item
Сигнал «АРВ Вкл.ВВ или СВ» с выхода терминала рабочего ввода поступает на дискретный вход «Включение АВР» терминала СВ (или резервного ввода) и при наличии условий готовности АВР и контроля отсутствия напряжения на секции шин (если контроль предусмотрен уставкой) с выходного реле терминала формируется команда в привод на включение СВ (или выключателя резервного ввода). 

\item
После включения выключателя формируется команда сброса готовности АВР. Сброс готовности АВР также формируется при поступлении сигналов из логики КСВ или выводом функции АВР из работы. 

\item
Поясняющие схемы организации цепей АВР и ВНР показаны на рисунках \ref{appavr:img1} и \ref{appavr:img2} приложения \ref{app:avr}. 

\item
Ввод АВР в работу на этапе параметрирования устройства осуществляется программным переключателем SGF1 «Ввод\_функции». Информация о введенной функции отображается выходным сигналом «АВР: Ввод».

\item
Функция может быть выведена из работы оперативно путем активации сигнала «ОВ АВР». Выведенное состояние функции характеризуется активным состоянием сигнала «АВР: Оперативный вывод».

\item
Алгоритм работы функции выполнен в соответствии с рисунком \ref{avr:img1}. В таблице \ref{avr:tbl1} приведены параметры, необходимые для настройки функции.

\vspace{1mm}
\begin{figure}[H]
\centering
\includegraphics[width=0.85\textwidth,height=0.85\textheight,keepaspectratio]{\fbpath/308. АВР/_latex/img/avr.pdf}
\captionof{figure}{Функциональная схема логики АВР}\label{avr:img1}
\end{figure}

\end{enumerate}

%%%%%%%%%%%%%%%%%%%%%%%%%%%%%%%%%%%%%%%%%%%%%%%%%%%%%%%%%%%%%%%%%%%%%% ЗАПРЕТ АВР %%%%%%%%%%%%%%%%%%%%%%%%%%%%%%%%%%%%%%%%%%%%%%%%%%
\item Формирование запрета АВР (ЗАВР)
\begin{enumerate}

\item
В устройстве запрет автоматического ввода резерва представлен функцией «ЗАВР». Запрет АВР позволяет избежать непреднамеренного включения оборудования, что может привести к повреждению или нарушению безопасности.

\item
Запрет АВР активируется, когда введена функция АВР и поступил один из сигналов как из логики, так и на дискретные входы устройства.

\item
Запрет АВР от устройства ТН, функции сигнализации замыкания на землю (СЗЗ) первой или второй секции шин контролируется уставкой SGF3. Возможен вывод действия сигналов «КСВ: В самопр. отключен»,  «КСВ: Неисправность В» на запрет АВР программными переключателями SGF1 и SGF2 соответственно.

\item
Алгоритм работы функции выполнен в соответствии с рисунком \ref{zavr:img1}. В таблице \ref{zavr:tbl1} приведены параметры, необходимые для настройки функции.

\vspace{3mm}
\begin{figure}[H]
\centering

\includegraphics[width=0.85\textwidth,height=0.85\textheight,keepaspectratio]{\fbpath/308. АВР/_latex/img/zavr.pdf}
\captionof{figure}{Функциональная схема логики ЗАВР}\label{zavr:img1}
\end{figure}

\end{enumerate}

\end{enumerate}
